\makeabstract
    {
    Kinetics of Solid-State Polymerization of H6TNAP
    }
    {
    Zen Loo\textsuperscript{1}, Emma Abell\textsuperscript{1}, Ren A. Wiscons\textsuperscript{1}
    }
    {
    \textsuperscript{1} Department of Chemistry, Amherst College, Amherst, MA
    }
    {
    11,11,12,12-Tetracyanonaphtho-2,6-quinodimethane, which we abbreviate as H6TNAP, is one of few known molecules to polymerize in the solid state. Unlike polymerization in the liquid or gas state, solid state polymerization provides a high degree of control over polymerization reactions. These controlled polymer syntheses have applications in materials science, specifically for producing copolymers with the capability to stabilize donor-acceptor interfaces in electronic materials, like semiconductors and photovolatics. However, due to the rarity of solid-state polymerization in organic materials, much is still unknown about their kinetics. Our project seeks to explore the kinetics of solid-state polymerization in an organic crystalline substance, focusing on how the polymerization propagates and the factors that explain why.
    }
    {
    H6TNAP Polymerization occurs under conditions of high heat and UV-Vis light. We created these conditions using a hot plate set between 200-220{\textordmasculine}C and a photoreactor, 3D printed to keep an LED light {\textasciitilde}2 cm away from the sample. We created samples on 40 mm glass cover slips, isolating 5-7 H6TNAP crystals on the center of each and sealing with another cover slip kept 0.06mm apart using one layer of Kapton tape. Because the H6TNAP monomer is red and its polymer is colorless, we collected data using a brightfield microscope, which accurately indicated the polymerization using average intensity measurements. Images were captured every 5 minutes and analyzed using Zen Imaging software to create time series for each trial.
    }
    {
    We observed both linear and quadratic average intensity time series. Linear time series were most often observed in data sets with less time points, while quadratic time series were observed in longer data sets, like those with larger crystals or lower temperature trials. Via image analysis, we observed that perimeter to area ratio increased over time and that core intensity increased over time as well.
    }
    {
    Given that shorter time series yielded linear data, we hypothesized that the standard polymerization kinetics of our organic crystals are linear by area. Under this assumption, our image analysis observations allowed us to identify factors that explained the quadratic average intensity time series, namely an increasing perimeter:area ratio and 3D polymerization, which we hypothesized to cause slight increases in polymerization rate over time in larger crystals or lower temperature trials. Currently, we do not have adequate data to necessarily prove our hypotheses. Our next steps are to improve methods as to run trials more efficiently and with less miscellaneous error.
    }

    \newpage
    